\documentclass[12pt,landscape]{exam}
\usepackage[margin=.3in]{geometry}
\setlength{\columnsep}{.6in}
\usepackage{multicol}

\begin{document}
\pagestyle{empty}

\newcommand{\mycontent}{
  \section*{Newton's First Law}
  \begin{tabular}{p{1.8in}|p{2.7in}}
    An object maintains its state of rest or uniform motion in a straight line unless acted on by a net force. &
    \vspace{1em}
    \raggedright
    An object maintains its \fillin[momentum][9em] unless acted on by an \fillin[outside impulse][11em]\\[1em]
  \end{tabular}
  \vspace{1em}
  \hrule

  \section*{Newton's Second Law}
  \begin{tabular}{p{1.8in}|p{2.7in}}
    Acceleration is inversely proportional to mass, directly proportional to net force, and in the same direction as the net force. & \\
    &\\[1em]
    \centering $F_{NET}=ma$ & \\
    &\\
    & \fillin[impulse][12em] is equal to \fillin[change in momentum][2.7in] \\[1em]
  \end{tabular}
  \vspace{1em}
  \hrule


  \section*{Newton's Third Law}
  \begin{tabular}{p{1.8in}|p{2.7in}}
    When one object apples a force on a second object, the second object applies an equal and opposite force on the first object. & \\
    &\\[1em]
    \centering $F_A=-F_B$ & \\
    &\\[3em]
    & Momentum is \fillin[conserved][10.4em]
  \end{tabular}
}

\begin{multicols}{2}

  \mycontent

  \columnbreak

  \mycontent

\end{multicols}

\end{document}